\section{Introduction}
\label{sec:intro}

Operational decision support for outdoor activities rests on \emph{suitability
scoring}: a mapping from the environmental state---wind speed, wave height,
thermal lift, snow stability---to a bounded score or a go/no-go verdict. In
practice these mappings are expert-defined curves, one per activity and per
physical dimension, calibrated by domain specialists and, increasingly, refined
against observed outcomes. A kitesurfing wind curve peaks near a sweet spot and
falls off toward calm and storm; a ski-touring stability curve gates on
avalanche risk. The score is then a composition of such curves.

A single suitability curve, however, silently conflates two quantities that a
decision-maker needs to keep apart. The first is the \emph{intrinsic difficulty}
of a condition: how objectively demanding a given wind or sea state is for the
activity, independent of who is attempting it. The second is the \emph{skill} of
the person facing that condition. When a curve is fit to behavioural outcomes---
did the session go well, did the rider go out at all---it absorbs both: the
apparent ``suitability'' of $22$ knots at a given spot reflects not only how hard
$22$ knots is, but also the skill distribution of whoever happens to ride there.
Two spots with identical fitted suitability curves can therefore have radically
different intrinsic difficulty, masked by different clienteles. No existing
observation network---marine, atmospheric, cryospheric---measures the difficulty
term in isolation, because none carries the human-decision layer needed to
identify it.

\paragraph{The gap.} What is missing is a way to \emph{separate} intrinsic
condition difficulty from population skill using the outcome data that decision
platforms already collect, and to do so with an identified, interpretable model
rather than a black box. The methodological obstacle is that conditions are not
discrete test items: wind speed is a continuous axis, not a yes/no question, so
the classical psychometric machinery for separating item difficulty from person
ability does not apply off the shelf.

\paragraph{Contributions.} This paper makes three contributions.
\begin{enumerate}
  \item \textbf{Continuous-item IRT for environmental conditions.} We extend
    explanatory item response theory---which already lets item difficulty depend
    on covariates \citep{fischer1973,deboeck2004}---from discrete item features to
    a \emph{continuous physical driver}, modelling difficulty $\delta(x,s)$ as a
    smooth, physically anchored function of a metric $x$ at site $s$
    (Section~\ref{sec:model}). The item bank is the weather.
  \item \textbf{Identifiability and synthetic recovery.} We give the conditions
    under which rider skill and condition difficulty separate---a connected
    rider${\times}$condition incidence graph plus location/scale anchoring---and a
    deterministic marginal-maximum-likelihood estimator with a documented
    single-curve fallback when identification fails (Section~\ref{sec:ident}). We
    validate by synthetic recovery: on a reference cohort the model recovers
    latent skill at correlation $0.96$ and locates the difficulty minimum within
    $\pm 3$ units of ground truth (Section~\ref{sec:recovery}).
  \item \textbf{The difficulty atlas.} The recovered difficulty function defines a
    measurable, site-level construct---an \emph{intrinsic difficulty atlas}---that
    is anonymous, joinable to existing meteorological archives, and captures a
    quantity those archives do not (Section~\ref{sec:atlas}).
\end{enumerate}

\paragraph{Scope.} This is a methodology paper. Its validation is a synthetic
recovery study: cohorts are generated with known ground-truth skill and
difficulty, and the estimator is shown to recover them. We make \emph{no}
empirical claim on real riders; no proprietary observations are used anywhere,
and every number in the paper reproduces from a single command on synthetic data
(Section~\ref{sec:recovery}). Empirical validation on field cohorts is future
work (Section~\ref{sec:discussion}).
